\documentclass[UTF8,11pt]{ctexart}
\usepackage[a4paper,margin=24mm]{geometry}
\usepackage{amsmath,amssymb,amsthm,mathtools,mathrsfs}
\usepackage[all]{xy}
\usepackage{xurl,hyperref,enumitem}
\usepackage{tikz}
\usetikzlibrary{arrows.meta,cd,decorations.markings,calc,patterns}
\hypersetup{hidelinks,breaklinks=true}
\setlength{\emergencystretch}{3em}
\allowdisplaybreaks
\newtheorem*{theorem}{Theorem}
\newtheorem*{thm}{Theorem}
\newtheorem*{lemma}{Lemma}
\newtheorem*{proposition}{Proposition}
\newtheorem*{prop}{Proposition}
\newtheorem*{corollary}{Corollary}
\newtheorem*{cor}{Corollary}
\newtheorem*{claim}{Claim}
\theoremstyle{definition}
\newtheorem*{definition}{Definition}
\newtheorem*{defn}{Definition}
\newtheorem*{remark}{Remark}
\newtheorem*{example}{Example}
\newcommand{\R}{\mathbb R}
\newcommand{\C}{\mathbb C}
\newcommand{\Z}{\mathbb Z}
\newcommand{\Q}{\mathbb Q}
\newcommand{\N}{\mathbb N}
\newcommand{\F}{\mathbb F}
\newcommand{\D}{\mathbb D}
\newcommand{\bB}{\mathbb B}
\newcommand{\bC}{\mathbb C}
\newcommand{\bR}{\mathbb R}
\newcommand{\bZ}{\mathbb Z}
\newcommand{\bN}{\mathbb N}
\newcommand{\bQ}{\mathbb Q}
\newcommand{\bD}{\mathbb D}
\newcommand{\bF}{\mathbb F}
\newcommand{\CP}{\mathbb{CP}}
\newcommand{\RP}{\mathbb{RP}}
\newcommand{\sO}{\mathscr O}
\newcommand{\sI}{\mathscr I}
\newcommand{\sF}{\mathscr F}
\newcommand{\sG}{\mathscr G}
\newcommand{\sH}{\mathscr H}
\newcommand{\sR}{\mathscr R}
\newcommand{\sM}{\mathscr M}
\newcommand{\sV}{\mathscr V}
\newcommand{\sU}{\mathscr U}
\DeclareMathOperator{\Hom}{Hom}
\DeclareMathOperator{\Ext}{Ext}
\DeclareMathOperator{\Tor}{Tor}
\DeclareMathOperator{\Spec}{Spec}
\DeclareMathOperator{\Proj}{Proj}
\DeclareMathOperator{\supp}{supp}
\DeclareMathOperator{\im}{im}
\DeclareMathOperator{\id}{id}
\DeclareMathOperator{\Tr}{Tr}
\DeclareMathOperator{\tr}{tr}
\DeclareMathOperator{\rank}{rank}
\DeclareMathOperator{\ord}{ord}
\DeclareMathOperator{\dist}{dist}
\DeclareMathOperator{\End}{End}
\DeclareMathOperator{\Aut}{Aut}
\DeclareMathOperator{\Gal}{Gal}
\DeclareMathOperator{\vol}{vol}
\DeclareMathOperator{\Cl}{Cl}
\DeclareMathOperator{\coker}{coker}
\DeclareMathOperator{\codim}{codim}
\DeclareMathOperator{\divergence}{div}
\DeclareMathOperator{\Ric}{Ric}
\DeclareMathOperator{\grad}{grad}
\DeclareMathOperator{\Hess}{Hess}
\newcommand{\abs}[1]{\left\lvert#1\right\rvert}
\newcommand{\sabs}[1]{\lvert#1\rvert}
\newcommand{\babs}[1]{\bigl\lvert#1\bigr\rvert}
\newcommand{\Babs}[1]{\Bigl\lvert#1\Bigr\rvert}
\newcommand{\bbabs}[1]{\biggl\lvert#1\biggr\rvert}
\newcommand{\BBabs}[1]{\Biggl\lvert#1\Biggr\rvert}
\newcommand{\norm}[1]{\left\lVert#1\right\rVert}
\newcommand{\snorm}[1]{\lVert#1\rVert}
\newcommand{\bnorm}[1]{\bigl\lVert#1\bigr\rVert}
\newcommand{\Bnorm}[1]{\Bigl\lVert#1\Bigr\rVert}
\newcommand{\bbnorm}[1]{\biggl\lVert#1\biggr\rVert}
\newcommand{\BBnorm}[1]{\Biggl\lVert#1\Biggr\rVert}
\newcommand{\widebar}[1]{\overline{#1}}
\newcommand{\nicefrac}[2]{\frac{#1}{#2}}
\newcommand{\linnprod}[2]{\langle#1,#2\rangle}
\newcommand{\myindex}[1]{#1}
\newcommand{\glsadd}[1]{}
\newcommand{\avoidbreak}{}
\newcommand{\sourceref}[1]{\textnormal{[原文：\nolinkurl{#1}]}}
\newcommand{\thmref}[1]{\sourceref{#1}}
\newcommand{\lemmaref}[1]{\sourceref{#1}}
\newcommand{\propref}[1]{\sourceref{#1}}
\newcommand{\corref}[1]{\sourceref{#1}}
\newcommand{\defnref}[1]{\sourceref{#1}}
\newcommand{\exampleref}[1]{\sourceref{#1}}
\newcommand{\exerciseref}[1]{\sourceref{#1}}
\newcommand{\figureref}[1]{\sourceref{#1}}
\newcommand{\sectionref}[1]{\sourceref{#1}}
\newcommand{\appendixref}[1]{\sourceref{#1}}
\newcommand{\stacksref}[2]{\href{https://stacks.math.columbia.edu/tag/#1}{#2 [#1]}}


\newcommand{\E}{\mathbb E}
\newcommand{\Prob}{\mathbb P}
\newcommand{\Reff}{\mathcal R_{\mathrm{eff}}}
\newcommand{\eps}{\varepsilon}
\DeclareMathOperator{\diam}{diam}
\DeclareMathOperator{\Var}{Var}
\DeclareMathOperator{\Crit}{Crit}
\newcommand{\dd}{\,\mathrm d}

\begin{document}
\section*{060\quad 代数闭域上函数域的Tsen性质}
\subsection*{来源与计数目标}
The Stacks Project Authors, \emph{The Stacks Project}, Theorem 59.67.10 (Tag 03RD)，Definition 59.67.5、Lemma 59.67.7；补列 Proposition 33.45.13 (0BJ8) 和 Lemma 33.34.2 (0B2Q)。使用2026-09-28访问的在线正文。\par\url{https://stacks.math.columbia.edu/tag/0A2M}\par\url{https://stacks.math.columbia.edu/tag/0BJ8}\par\url{https://stacks.math.columbia.edu/tag/0B2Q}
\par\url{https://stacks.math.columbia.edu/tag/03RD}
\par\textbf{本文件只计一个问题：}证明 $r$ 维簇的函数域为 $C_r$；后面的 Brauer 群推论及前置引理不另外计数。
\subsection*{作者命题与人类证明}
\begin{definition}[59.67.5]
A field $K$ is called $C_r$ if for every $0<d^r<n$ and every $f\in K[T_1,\ldots,T_n]$ homogeneous of degree $d$, there exist $\alpha=(\alpha_1,\ldots,\alpha_n)$, $\alpha_i\in K$ not all zero, such that $f(\alpha)=0$. Such an $\alpha$ is called a nontrivial solution of $f$.
\end{definition}
\begin{lemma}[59.67.7]
Let $k$ be an algebraically closed field. Let $f_1,\ldots,f_s\in k[T_1,\ldots,T_n]$ be homogeneous polynomials of degree $d_1,\ldots,d_s$ with $d_i>0$. If $s<n$, then $f_1=\ldots=f_s=0$ have a common nontrivial solution.
\end{lemma}
\begin{proof}
This follows from dimension theory, for example in the form of Varieties, Lemma 33.34.2 applied $s-1$ times.
\end{proof}
\begin{theorem}[59.67.10, Tsen's theorem]
The function field of a variety of dimension $r$ over an algebraically closed field $k$ is $C_r$.
\end{theorem}
\begin{proof}
For projective space one can show directly that the field $k(x_1,\ldots,x_r)$ is $C_r$ (exercise).

General case. Without loss of generality, we may assume $X$ to be projective. Let $f\in k(X)[T_1,\ldots,T_n]_d$ with $0<d^r<n$. Say the coefficients of $f$ are in $\Gamma(X,\mathcal O_X(H))$ for some ample $H\subset X$. Let $\boldsymbol\alpha=(\alpha_1,\ldots,\alpha_n)$ with $\alpha_i\in\Gamma(X,\mathcal O_X(eH))$. Then $f(\boldsymbol\alpha)\in\Gamma(X,\mathcal O_X((de+1)H))$. Consider the system of equations $f(\boldsymbol\alpha)=0$. Then by asymptotic Riemann--Roch (Varieties, Proposition 33.45.13) there exists a $c>0$ such that
\begin{itemize}
\item the number of variables is $n\dim_k\Gamma(X,\mathcal O_X(eH))\sim ne^rc$, and
\item the number of equations is $\dim_k\Gamma(X,\mathcal O_X((de+1)H))\sim(de+1)^rc$.
\end{itemize}
Since $n>d^r$, there are more variables than equations. The equations are homogeneous hence there is a solution by Lemma 59.67.7.
\end{proof}
\paragraph{Quoted standard prerequisite: asymptotic Riemann--Roch.}
\begin{proposition}[33.45.13, Tag 0BJ8]
Let $k$ be a field. Let $X$ be a proper scheme over $k$ of dimension $d$. Let $\mathcal L$ be an ample invertible $\mathcal O_X$-module. Then
\[\dim_k\Gamma(X,\mathcal L^{\otimes n})\sim cn^d+l.o.t.\]
where $c=\deg_{\mathcal L}(X)/d!$ is a positive constant.
\end{proposition}
\begin{proof}
This follows from the definitions, Lemma 33.45.9, and the vanishing of higher cohomology in Cohomology of Schemes, Lemma 30.17.1.
\end{proof}
\paragraph{Dimension-theoretic input used in Lemma 59.67.7.}
\begin{lemma}[33.34.2, Tag 0B2Q]
Let $k$ be a field and $n\geq0$. Let $X,Y\subset\mathbf A^n_k$ be closed subsets. Assume that $X$ and $Y$ are equidimensional, $\dim(X)=r$ and $\dim(Y)=s$. Then every irreducible component of $X\cap Y$ has dimension $\geq r+s-n$.
\end{lemma}
\begin{proof}
Consider the closed subscheme $X\times Y\subset\mathbf A^{2n}_k$ where we use coordinates $x_1,\ldots,x_n,y_1,\ldots,y_n$. Then $X\cap Y=X\times Y\cap V(x_1-y_1,\ldots,x_n-y_n)$. Let $t\in X\cap Y\subset X\times Y$ be a closed point. By Lemma 33.20.5 we have $\dim_t(X\times Y)=\dim(X)+\dim(Y)$. Thus $\dim(\mathcal O_{X\times Y,t})=r+s$ by Lemma 33.20.3. By Algebra, Lemma 10.60.13 we conclude that
\[\dim(\mathcal O_{X\cap Y,t})=\dim(\mathcal O_{X\times Y,t}/(x_1-y_1,\ldots,x_n-y_n))\geq r+s-n.\]
This implies the result by Lemma 33.20.3.
\end{proof}
\subsection*{整理说明（非作者正文）}
主证明完整保留；原文首句的射影空间直接证明留作练习，但随后的 General case 独立覆盖一般簇，并不调用该练习作为论据。后列渐近 Riemann--Roch 的作者陈述与证明、交维数论证仅供读者识别被引用工具。明确允许的标准先修为射影模型、丰沛除子的截面与有理函数、Hilbert 多项式/Serre 消失所给截面维数渐近，以及仿射维数理论；本题的核心工作是为未知解选择增长的截面空间，把函数域上的单个方程转为基域上多个齐次方程，再比较变量数与方程数。没有把 $C_r$ 性质或 Tsen 定理作为先修。正文为网页数学 TeX 摘录，未取得固定网页构建 commit；不编造版本。
\subsection*{可修改的简短标注（非作者正文）}
$c$：函数域上齐次方程非平凡解的存在性

$a$：丰沛除子；截面空间；渐近Riemann--Roch；齐次方程组；代数簇交维数

\texttt{cross\_theory\_status = yes}。将函数域方程的未知量限制到截面空间，把几何维数控制转成齐次方程组的变量数优势，再返回函数域中的非零解。位置：59.67.10的General case。
标签为非穷尽、可修改初稿。
\subsection*{许可与修改记录}
原作者及作品见来源栏；来源采用 GNU Free Documentation License 1.2 或后续版本。本题证摘录和附加整理沿用该许可。2026-09-28：增加题号、来源定位、独立排版、整理说明和初稿标注；数学内容的取材范围及排版变化见整理说明。
\end{document}
